\section{Simulaciones: dímero de dos esferas}
\label{sec:dimero}

En esta sección se estudia el sistema de dispersión múltiple más simple: un
dímero formado por dos esferas de oro separadas por un hueco de
anchura $d$ (gap). El eje que une los centros de las dos esferas está orientado
a lo largo del eje $x$ (configuración E--W), con la onda incidente
propagándose a lo largo del eje $z$. Esta orientación se escoge porque el
campo cercano de una esfera individual se concentra en este eje, lo que
maximiza el acoplamiento entre las dos partículas.

Se consideran tres radios que cubren regímenes electromagnéticos
cualitativamente distintos, resumidos en la tabla~\ref{tab:dimero_regimenes}.

\begin{table}[H]
  \centering
  \begin{tabular}{ccc}
    \hline
    Radio $a$ & Rango de $x$ ($\lambda\in[400,800]\,\text{nm}$) & Régimen \\
    \hline
    25\,nm  & $0.20$ -- $0.39$ & Dipolar (Rayleigh) \\
    85\,nm  & $0.67$ -- $1.34$ & Transición (cuadrupolar) \\
    600\,nm & $4.71$ -- $9.42$ & Mie avanzado \\
    \hline
  \end{tabular}
  \caption{Radios simulados en el dímero E--W y régimen electromagnético
           correspondiente en el rango visible.}
  \label{tab:dimero_regimenes}
\end{table}

Para cada radio se estudia el efecto del gap variando $d$ desde distancias
muy pequeñas (acoplamiento fuerte) hasta distancias grandes 
(acoplamiento débil). En cada caso se presentan: las eficiencias
espectrales del conjunto y el campo cercano $|\mathbf{E}|^2$ evaluado en el
plano $xz$ a la longitud de onda de resonancia del dímero.


% =========================================================================
\subsection{Esfera individual: referencia}
\label{subsec:dimero_ref}
% =========================================================================

Antes de analizar los dímeros se recogen los resultados de la esfera individual
para cada radio, que servirán como referencia para cuantificar el efecto del
acoplamiento.

\begin{figure}[H]
  \centering
  \begin{subfigure}[b]{0.32\textwidth}
    \centering
    \includegraphics[width=\textwidth]{../../Simulaciones/Definitivo/simulacion-de-referencia/1e/ew/a25/sw/plots/Eficiencias_-_Unpol.png}
    \caption{$a=25\,\text{nm}$.}
  \end{subfigure}
  \hfill
  \begin{subfigure}[b]{0.32\textwidth}
    \centering
    \includegraphics[width=\textwidth]{../../Simulaciones/Definitivo/simulacion-de-referencia/1e/ew/a85/sw/plots/Eficiencias_-_Unpol.png}
    \caption{$a=85\,\text{nm}$.}
  \end{subfigure}
  \hfill
  \begin{subfigure}[b]{0.32\textwidth}
    \centering
    \includegraphics[width=\textwidth]{../../Simulaciones/Definitivo/simulacion-de-referencia/1e/ew/a600/sw/plots/Eficiencias_-_Unpol.png}
    \caption{$a=600\,\text{nm}$.}
  \end{subfigure}
  \caption{Eficiencias espectrales de la esfera individual de Au para los tres
            radios estudiados. Para $a=25\,\text{nm}$ la respuesta óptica está
            dominada por el modo dipolar con $Q_\text{abs}$ dominante ($x\ll1$).
            Para $a=85\,\text{nm}$ el espectro presenta un pico pronunciado en
            torno a $\lambda\approx570\,\text{nm}$ donde la contribución
            cuadrupolar es relevante ($x\sim1$). Para $a=600\,\text{nm}$ las
            eficiencias varían suavemente con contribuciones multipolares de
            orden superior, siendo la dispersión dominante en todo el rango
            visible ($x\gg1$).}
  \label{fig:ref_eff}
\end{figure}

\begin{figure}[H]
  \centering
  \begin{subfigure}[b]{0.32\textwidth}
    \centering
    \includegraphics[width=\textwidth]{../../Simulaciones/Definitivo/simulacion-de-referencia/1e/ew/a25/wl520/plots/Figure_1.png}
    \caption{$a=25\,\text{nm}$, $\lambda=520\,\text{nm}$.
             Máximo: $11.90\,|\mathbf{E}_0|^2$.}
  \end{subfigure}
  \hfill
  \begin{subfigure}[b]{0.32\textwidth}
    \centering
    \includegraphics[width=\textwidth]{../../Simulaciones/Definitivo/simulacion-de-referencia/1e/ew/a85/wl570/plots/Figure_1.png}
    \caption{$a=85\,\text{nm}$, $\lambda=570\,\text{nm}$.
             Máximo: $12.75\,|\mathbf{E}_0|^2$.}
  \end{subfigure}
  \hfill
  \begin{subfigure}[b]{0.32\textwidth}
    \centering
    \includegraphics[width=\textwidth]{../../Simulaciones/Definitivo/simulacion-de-referencia/1e/ew/a600/wl560/plots/Figure_1.png}
    \caption{$a=600\,\text{nm}$, $\lambda=560\,\text{nm}$.
             Máximo: $6.55\,|\mathbf{E}_0|^2$.}
  \end{subfigure}
  \caption{Campo cercano $|\mathbf{E}|^2$ de la esfera individual de Au para
           los tres radios. 
           %evaluado a la longitud de onda de resonancia de cada uno. 
           Para $a=25\,\text{nm}$ la intensificación se localiza en
           los polos de la esfera perpendiculares a la dirección de
           propagación, característico del modo dipolar ($x\ll1$). Para
           $a=85\,\text{nm}$ la concentración superficial aparece ligeramente
           desplazada hacia la cara de iluminación, reflejo de la contribución
           de órdenes multipolares superiores ($x\sim1$). Para
           $a=600\,\text{nm}$ el campo exhibe patrones de interferencia
           característicos del régimen de Mie avanzado.}
  \label{fig:ref_nf}
\end{figure}


% =========================================================================
\subsection{Dímero en régimen dipolar: $a = 25\,\text{nm}$}
\label{subsec:dimero_a25}
% =========================================================================

El radio $a=25\,\text{nm}$ sitúa a cada esfera en el límite superior del
régimen cuasidipolar ($x\in[0.20,0.39]$), donde la respuesta óptica está
dominada por el término dipolar eléctrico $a_1$. Se han simulado tres
separaciones: $d=5$, $10$ y $50\,\text{nm}$.

\subsubsection*{Eficiencias espectrales}

\begin{figure}[H]
  \centering
  \begin{subfigure}[b]{0.32\textwidth}
    \centering
    \includegraphics[width=\textwidth]{../../Simulaciones/Definitivo/simulacion-de-referencia/2e/ew/a25/sep5/sw/plots/Eficiencias_-_Unpol.png}
    \caption{$d=5\,\text{nm}$.}
    \label{fig:a25_eff_g5}
  \end{subfigure}
  \hfill
  \begin{subfigure}[b]{0.32\textwidth}
    \centering
    \includegraphics[width=\textwidth]{../../Simulaciones/Definitivo/simulacion-de-referencia/2e/ew/a25/sep10/sw/plots/Eficiencias_-_Unpol.png}
    \caption{$d=10\,\text{nm}$.}
    \label{fig:a25_eff_g10}
  \end{subfigure}
  \hfill
  \begin{subfigure}[b]{0.32\textwidth}
    \centering
    \includegraphics[width=\textwidth]{../../Simulaciones/Definitivo/simulacion-de-referencia/2e/ew/a25/sep50/sw/plots/Eficiencias_-_Unpol.png}
    \caption{$d=50\,\text{nm}$.}
    \label{fig:a25_eff_g50}
  \end{subfigure}
  \caption{Eficiencias espectrales del dímero con $a=25\,\text{nm}$ para
           tres separaciones. La resonancia permanece en torno a
           $\lambda\approx520\,\text{nm}$ en los tres casos; el acoplamiento
           incrementa la amplitud del pico sin desplazarlo espectralmente.
           La absorción sigue siendo el canal dominante de extinción para
           todos los gaps estudiados. A medida que aumenta $d$, la amplitud
           del pico disminuye progresivamente, convergiendo hacia los valores
           de la esfera individual.}
  \label{fig:a25_eff}
\end{figure}

\subsubsection*{Campo cercano}

\begin{figure}[H]
  \centering
  \begin{subfigure}[b]{0.32\textwidth}
    \centering
    \includegraphics[width=\textwidth]{../../Simulaciones/Definitivo/simulacion-de-referencia/2e/ew/a25/sep5/wl530/plots/Figure_1.png}
    \caption{$d=5\,\text{nm}$, $\lambda=530\,\text{nm}$.
             Máximo: $147.35\,|\mathbf{E}_0|^2$.}
    \label{fig:a25_nf_g5}
  \end{subfigure}
  \hfill
  \begin{subfigure}[b]{0.32\textwidth}
    \centering
    \includegraphics[width=\textwidth]{../../Simulaciones/Definitivo/simulacion-de-referencia/2e/ew/a25/sep10/wl525/plots/Figure_1.png}
    \caption{$d=10\,\text{nm}$, $\lambda=525\,\text{nm}$.
             Máximo: $61.58\,|\mathbf{E}_0|^2$.}
    \label{fig:a25_nf_g10}
  \end{subfigure}
  \hfill
  \begin{subfigure}[b]{0.32\textwidth}
    \centering
    \includegraphics[width=\textwidth]{../../Simulaciones/Definitivo/simulacion-de-referencia/2e/ew/a25/sep50/wl520/plots/Figure_1.png}
    \caption{$d=50\,\text{nm}$, $\lambda=520\,\text{nm}$.
             Máximo: $14.67\,|\mathbf{E}_0|^2$.}
    \label{fig:a25_nf_g50}
  \end{subfigure}
  \caption{Campo cercano $|\mathbf{E}_\text{avg}|^2$ del dímero con
           $a=25\,\text{nm}$ evaluado a la longitud de onda de máxima
           $Q_\text{sca}$. La zona de concentración se localiza en el gap entre
           las dos esferas. Al reducir el gap de $50\,\text{nm}$ a
           $5\,\text{nm}$, la intensificación máxima pasa de
           $14.67\,|\mathbf{E}_0|^2$ a $147.35\,|\mathbf{E}_0|^2$,
           un aumento de un orden de magnitud respecto a la esfera
           individual ($11.90\,|\mathbf{E}_0|^2$).}
  \label{fig:a25_nf}
\end{figure}

El dímero de $a=25\,\text{nm}$ muestra un acoplamiento débil en el espectro de eficiencias, 
ya que apenas se observa desplazamiento de la resonancia. Sin embargo, en 
campo cercano el acoplamiento en el gap es intenso para separaciones pequeñas.


% =========================================================================
\subsection{Dímero en régimen de transición: $a = 85\,\text{nm}$}
\label{subsec:dimero_a85}
% =========================================================================

Con $a=85\,\text{nm}$ el parámetro de tamaño entra en el rango $x\in[0.67,1.34]$,
donde los modos cuadrupolares comienzan a contribuir de forma apreciable junto
al dipolar. Se han simulado cuatro separaciones: $d=4$, $10$, $50$ y
$150\,\text{nm}$, que barren desde el acoplamiento muy fuerte hasta el límite
de dispersión independiente.

\subsubsection*{Eficiencias espectrales}

\begin{figure}[H]
  \centering
  \begin{subfigure}[b]{0.45\textwidth}
    \centering
    \includegraphics[width=\textwidth]{../../Simulaciones/Definitivo/simulacion-de-referencia/2e/ew/sep4/sw/plots/Eficiencias_-_Unpol.png}
    \caption{$d=4\,\text{nm}$. Resonancia en $\lambda\approx720\,\text{nm}$.}
    \label{fig:a85_eff_g4}
  \end{subfigure}
  \hfill
  \begin{subfigure}[b]{0.45\textwidth}
    \centering
    \includegraphics[width=\textwidth]{../../Simulaciones/Definitivo/simulacion-de-referencia/2e/ew/sep10/sw/plots/Eficiencias_-_Unpol.png}
    \caption{$d=10\,\text{nm}$. Resonancia en $\lambda\approx700\,\text{nm}$.}
    \label{fig:a85_eff_g10}
  \end{subfigure}

  \vspace{0.5em}

  \begin{subfigure}[b]{0.45\textwidth}
    \centering
    \includegraphics[width=\textwidth]{../../Simulaciones/Definitivo/simulacion-de-referencia/2e/ew/sep50/sw/plots/Eficiencias_-_Unpol.png}
    \caption{$d=50\,\text{nm}$. Resonancia en $\lambda\approx600\,\text{nm}$.}
    \label{fig:a85_eff_g50}
  \end{subfigure}
  \hfill
  \begin{subfigure}[b]{0.45\textwidth}
    \centering
    \includegraphics[width=\textwidth]{../../Simulaciones/Definitivo/simulacion-de-referencia/2e/ew/sep150/sw/plots/Eficiencias_-_Unpol.png}
    \caption{$d=150\,\text{nm}$. Resonancia en $\lambda\approx560\,\text{nm}$.}
    \label{fig:a85_eff_g150}
  \end{subfigure}

  \caption{Eficiencias espectrales del dímero con $a=85\,\text{nm}$ para
           $d=4$, $10$, $50$ y $150\,\text{nm}$. Al reducir el gap, la
           resonancia de $Q_\text{sca}$ se desplaza hacia el rojo: de
           $\sim560\,\text{nm}$ (esfera individual) hasta $\sim720\,\text{nm}$
           para $d=4\,\text{nm}$, un desplazamiento de más de $160\,\text{nm}$.
           Simultáneamente, $Q_\text{abs}$ se reduce notablemente en la resonancia,
           indicando que la extinción queda dominada por la dispersión.
           Para $d=150\,\text{nm}$ el espectro recupera el perfil de la
           esfera individual.}
  \label{fig:a85_eff}
\end{figure}

\subsubsection*{Campo cercano}

\begin{figure}[H]
  \centering
  \begin{subfigure}[b]{0.45\textwidth}
    \centering
    \includegraphics[width=\textwidth]{../../Simulaciones/Definitivo/simulacion-de-referencia/2e/ew/sep4/wl720/plots/Figure_1_log.png}
    \caption{$d=4\,\text{nm}$, $\lambda=720\,\text{nm}$.
             Máximo: $174.31\,|\mathbf{E}_0|^2$.}
    \label{fig:a85_nf_g4}
  \end{subfigure}
  \hfill
  \begin{subfigure}[b]{0.45\textwidth}
    \centering
    \includegraphics[width=\textwidth]{../../Simulaciones/Definitivo/simulacion-de-referencia/2e/ew/sep10/wl700/plots/Figure_1_log.png}
    \caption{$d=10\,\text{nm}$, $\lambda=700\,\text{nm}$.
             Máximo: $106.03\,|\mathbf{E}_0|^2$.}
    \label{fig:a85_nf_g10}
  \end{subfigure}
\end{figure}

\begin{figure}[H]\ContinuedFloat
  \centering
  \begin{subfigure}[b]{0.45\textwidth}
    \centering
    \includegraphics[width=\textwidth]{../../Simulaciones/Definitivo/simulacion-de-referencia/2e/ew/sep50/wl600/plots/Figure_1.png}
    \caption{$d=50\,\text{nm}$, $\lambda=600\,\text{nm}$.
             Máximo: $22.68\,|\mathbf{E}_0|^2$.}
    \label{fig:a85_nf_g50}
  \end{subfigure}
  \hfill
  \begin{subfigure}[b]{0.45\textwidth}
    \centering
    \includegraphics[width=\textwidth]{../../Simulaciones/Definitivo/simulacion-de-referencia/2e/ew/sep150/wl560/plots/Figure_1.png}
    \caption{$d=150\,\text{nm}$, $\lambda=560\,\text{nm}$.
             Máximo: $10.77\,|\mathbf{E}_0|^2$, recuperando el campo
             de la esfera individual.}
    \label{fig:a85_nf_g150}
  \end{subfigure}
  \caption{Campo cercano $|\mathbf{E}_\text{avg}|^2$ del dímero con
           $a=85\,\text{nm}$ evaluado a la longitud de onda de máxima
           $Q_\text{sca}$. La zona de concentración se localiza en el gap y
           crece de $10.77$ (esfera individual, $d=150\,\text{nm}$) a
           $174.31\,|\mathbf{E}_0|^2$ ($d=4\,\text{nm}$).}
  \label{fig:a85_nf}
\end{figure}


El dímero opera en régimen de acoplamiento fuerte.
Para gaps pequeños el espectro muestra dos picos bien diferenciados: el
primero, a longitudes de onda más cortas, está dominado por la absorción;
el segundo, desplazado hacia el rojo, está dominado por la dispersión y
presenta mayor amplitud. La resonancia de este segundo pico se desplaza
más de $160\,\text{nm}$ respecto a la esfera individual, alcanzando
$\lambda\approx720\,\text{nm}$ para $d=4\,\text{nm}$. A medida que el
gap aumenta, ambos picos se acercan y acaban fusionándose: la absorción
recupera peso en la extinción total, aunque la dispersión sigue siendo
el canal dominante. Para $d=150\,\text{nm}$ los dos picos han convergido
y el espectro recupera prácticamente el perfil de la esfera individual,
indicando que el acoplamiento es ya despreciable.


% =========================================================================
\subsection{Dímero en régimen Mie avanzado: $a = 600\,\text{nm}$}
\label{subsec:dimero_a600}
% =========================================================================

Con $a=600\,\text{nm}$ el parámetro de tamaño alcanza valores
$x\in[4.71,9.42]$, muy superiores a la unidad. En este límite participan
múltiples órdenes multipolares y el espectro óptico pierde el carácter
resonante estrecho de los tamaños anteriores. Se han simulado tres
separaciones: $d=10$, $20$ y $60\,\text{nm}$.

\subsubsection*{Eficiencias espectrales}

\begin{figure}[H]
  \centering
  \begin{subfigure}[b]{0.32\textwidth}
    \centering
    \includegraphics[width=\textwidth]{../../Simulaciones/Definitivo/simulacion-de-referencia/2e/ew/a600/sep10/sw/plots/Eficiencias_-_Unpol.png}
    \caption{$d=10\,\text{nm}$.}
    \label{fig:a600_eff_g10}
  \end{subfigure}
  \hfill
  \begin{subfigure}[b]{0.32\textwidth}
    \centering
    \includegraphics[width=\textwidth]{../../Simulaciones/Definitivo/simulacion-de-referencia/2e/ew/a600/sep20/plots/Eficiencias_-_Unpol.png}
    \caption{$d=20\,\text{nm}$.}
    \label{fig:a600_eff_g20}
  \end{subfigure}
  \hfill
  \begin{subfigure}[b]{0.32\textwidth}
    \centering
    \includegraphics[width=\textwidth]{../../Simulaciones/Definitivo/simulacion-de-referencia/2e/ew/a600/sep60/plots/Eficiencias_-_Unpol.png}
    \caption{$d=60\,\text{nm}$.}
    \label{fig:a600_eff_g60}
  \end{subfigure}
  \caption{Eficiencias espectrales del dímero con $a=600\,\text{nm}$ para
           las tres separaciones estudiadas. Los espectros son prácticamente
           idénticos entre sí y al de la esfera individual: el gap no produce
           desplazamiento espectral ni modifica la forma del espectro.}
  \label{fig:a600_eff}
\end{figure}

\subsubsection*{Campo cercano}

\begin{figure}[H]
  \centering
  \begin{subfigure}[b]{0.32\textwidth}
    \centering
    \includegraphics[width=\textwidth]{../../Simulaciones/Definitivo/simulacion-de-referencia/2e/ew/a600/sep10/wl670/plots/Figure_1.png}
    \caption{$d=10\,\text{nm}$, $\lambda=670\,\text{nm}$.
              Máximo en gap: $15.02\,|\mathbf{E}_0|^2$.}
    \label{fig:a600_nf_g10}
  \end{subfigure}
  \hfill
  \begin{subfigure}[b]{0.32\textwidth}
    \centering
    \includegraphics[width=\textwidth]{../../Simulaciones/Definitivo/simulacion-de-referencia/2e/ew/a600/sep20/wl660/plots/Figure_1.png}
    \caption{$d=20\,\text{nm}$, $\lambda=660\,\text{nm}$.
             Máximo en gap: $11.52\,|\mathbf{E}_0|^2$.}
    \label{fig:a600_nf_g20}
  \end{subfigure}
  \hfill
  \begin{subfigure}[b]{0.32\textwidth}
    \centering
    \includegraphics[width=\textwidth]{../../Simulaciones/Definitivo/simulacion-de-referencia/2e/ew/a600/sep60/wl670/plots/Figure_1.png}
    \caption{$d=60\,\text{nm}$, $\lambda=670\,\text{nm}$.
             Máximo en gap: $6.44\,|\mathbf{E}_0|^2$.}
    \label{fig:a600_nf_g60}
  \end{subfigure}
  \caption{Campo cercano $|\mathbf{E}_\text{avg}|^2$ del dímero con
           $a=600\,\text{nm}$. El rasgo dominante no es la zona de
           concentración del gap sino los patrones de interferencia que
           rodean a las esferas.}
  \label{fig:a600_nf}
\end{figure}

En el régimen Mie avanzado el acoplamiento tiene un efecto limitado en campo
lejano ya que el espectro de eficiencias apenas varía con la separación.
En campo cercano la zona de concentración en el gap sigue presente, pero
los patrones de interferencia de cada esfera dominan la imagen. Este resultado
sugiere que, en este régimen, el dímero es menos eficaz como nanoantena óptica
que a tamaños menores.

% =========================================================================
\subsection{Dímero heterogéneo: $a_1 = 25\,\text{nm}$, $a_2 = 85\,\text{nm}$}
\label{subsec:dimero_hetero}
% =========================================================================

La combinación de una esfera pequeña ($a_1=25\,\text{nm}$, régimen dipolar) con
una esfera mediana ($a_2=85\,\text{nm}$, contribución cuadrupolar) permite explorar
el acoplamiento asimétrico entre dos modos de naturaleza distinta. Se han
simulado tres separaciones: $d=5$, $10$ y $50\,\text{nm}$.

\subsubsection*{Eficiencias espectrales}

\begin{figure}[H]
  \centering
  \begin{subfigure}[b]{0.32\textwidth}
    \centering
    \includegraphics[width=\textwidth]{../../Simulaciones/Definitivo/simulacion-de-referencia/2e/ew/a85a25/sep5/sw/plots/Eficiencias_-_Unpol.png}
    \caption{$d=5\,\text{nm}$.}
  \end{subfigure}
  \hfill
  \begin{subfigure}[b]{0.32\textwidth}
    \centering
    \includegraphics[width=\textwidth]{../../Simulaciones/Definitivo/simulacion-de-referencia/2e/ew/a85a25/sep10/sw/plots/Eficiencias_-_Unpol.png}
    \caption{$d=10\,\text{nm}$.}
  \end{subfigure}
  \hfill
  \begin{subfigure}[b]{0.32\textwidth}
    \centering
    \includegraphics[width=\textwidth]{../../Simulaciones/Definitivo/simulacion-de-referencia/2e/ew/a85a25/sep50/sw/plots/Eficiencias_-_Unpol.png}
    \caption{$d=50\,\text{nm}$.}
  \end{subfigure}
  \caption{Eficiencias espectrales del dímero heterogéneo
           ($a_1=25\,\text{nm}$, $a_2=85\,\text{nm}$) para tres separaciones.
           El espectro reproduce la forma de la esfera individual de
           $85\,\text{nm}$: pico de dispersión dominante en torno a
           $\lambda\approx560$--$580\,\text{nm}$ con absorción pequeña.
           La esfera mayor gobierna la respuesta óptica del conjunto y
           el gap no produce desplazamiento espectral apreciable.}
  \label{fig:hetero_eff}
\end{figure}

\subsubsection*{Campo cercano}

\begin{figure}[H]
  \centering
  \begin{subfigure}[b]{0.32\textwidth}
    \centering
    \includegraphics[width=\textwidth]{../../Simulaciones/Definitivo/simulacion-de-referencia/2e/ew/a85a25/sep5/wl590/plots/Figure_1.png}
    \caption{$d=5\,\text{nm}$, $\lambda=590\,\text{nm}$.
             Máximo: $186.16\,|\mathbf{E}_0|^2$.}
  \end{subfigure}
  \hfill
  \begin{subfigure}[b]{0.32\textwidth}
    \centering
    \includegraphics[width=\textwidth]{../../Simulaciones/Definitivo/simulacion-de-referencia/2e/ew/a85a25/sep10/wl585/plots/Figure_1.png}
    \caption{$d=10\,\text{nm}$, $\lambda=585\,\text{nm}$.
             Máximo: $117.51\,|\mathbf{E}_0|^2$.}
  \end{subfigure}
  \hfill
  \begin{subfigure}[b]{0.32\textwidth}
    \centering
    \includegraphics[width=\textwidth]{../../Simulaciones/Definitivo/simulacion-de-referencia/2e/ew/a85a25/sep50/wl580/plots/Figure_1.png}
    \caption{$d=50\,\text{nm}$, $\lambda=580\,\text{nm}$.
             Máximo: $20.97\,|\mathbf{E}_0|^2$.}
  \end{subfigure}
  \caption{Campo cercano $|\mathbf{E}_\text{avg}|^2$ del dímero heterogéneo
           evaluado a la longitud de onda de resonancia. La zona de
           concentración se localiza en el gap, entre la cara de la esfera
           grande y la esfera pequeña. Al reducir $d$ de $50$ a
           $5\,\text{nm}$, la intensificación crece de $20.97$ a
           $186.16\,|\mathbf{E}_0|^2$.}
  \label{fig:hetero_nf}
\end{figure}

En el dímero heterogéneo la respuesta óptica está gobernada por la esfera
mayor: el espectro de eficiencias reproduce el perfil de la partícula de
$85\,\text{nm}$ individual, sin desplazamiento espectral apreciable con el
gap. Sin embargo, la asimetría geométrica concentra el campo en el lado de
la esfera pequeña, alcanzando $186\,|\mathbf{E}_0|^2$ para $d=5\,\text{nm}$,
comparable a lo obtenido en el dímero homogéneo de $85\,\text{nm}$.


% % =========================================================================
% \subsection{Comparación entre regímenes}
% \label{subsec:dimero_comparacion}
% % =========================================================================

% Los tres tamaños estudiados ilustran tres comportamientos cualitativamente
% distintos del dímero al variar el gap, que se resumen en la
% tabla~\ref{tab:dimero_resumen}.

% \begin{table}[H]
%   \centering
%   \small
%   \begin{tabular}{lccc}
%     \hline
%     & $a=25\,\text{nm}$ & $a=85\,\text{nm}$ & $a=600\,\text{nm}$ \\
%     \hline
%     Desplazamiento espectral al cerrar gap  & Nulo          & Fuerte ($>160\,\text{nm}$) & Nulo \\
%     Naturaleza de la extinción en resonancia & Absorción     & Dispersión pura            & Dispersión \\
%     Hot spot máximo (gap mínimo)            & $147\times$   & $174\times$                & $\sim12\times$ \\
%     Sensibilidad al gap                     & Moderada      & Muy alta                   & Nula \\
%     \hline
%   \end{tabular}
%   \caption{Comparación entre los tres dímeros E--W de Au simulados.}
%   \label{tab:dimero_resumen}
% \end{table}

% La sensibilidad al gap resulta máxima para $a=85\,\text{nm}$, donde el
% acoplamiento produce tanto un desplazamiento espectral pronunciado como la
% mayor amplificación de campo cercano. El radio de $25\,\text{nm}$ ofrece
% también una amplificación elevada pero sin desplazamiento espectral, lo que
% puede ser ventajoso en aplicaciones que requieran amplificación selectiva a
% una longitud de onda fija. En cambio, para $a=600\,\text{nm}$ el acoplamiento
% es irrelevante en campo lejano y modesto en campo cercano: el régimen Mie
% avanzado hace que el tamaño de la partícula sea el factor determinante, no la
% geometría del dímero.